\scanpage{25}

\uline{Cas} $j=i-1$

\diagram{\begin{tikzcd}[ampersand replacement=\&,column sep=small,row sep=small]
\& \Psi^{i-1}\arrow[r,dash]\arrow[d,"n\,\mathrm{id}"'] \& \Phi^{(i)}\arrow[r]\arrow[d] \& \Psi^i\arrow[d,equal] \&\\
0\arrow[r,dash] \& \Psi^{i-1}\arrow[r,dash,"D"']\arrow[d,<->,decorate,decoration={snake,amplitude=.35mm,segment length=2mm}] \& \Phi(i,n)\arrow[r,dash]\arrow[d,"\alpha=?"] \& \Psi^i\arrow[r,dash]\arrow[d,<->,decorate,decoration={snake,amplitude=.35mm,segment length=2mm}] \& 0\\
0\arrow[r] \& \Psi^{i\deleted{\unsure{-1|+2}}-2}\arrow[r] \& \Phi(i-1,m)\arrow[r,dash,"u"] \& \Psi^{i-1\deleted{\unsure{+1}}}\arrow[r,dash] \& 0
\end{tikzcd}}

le composé $\Phi^i\longrightarrow\Phi(i,n)\mathrel{-}\Phi(i-1,m)\mathrel{-}\Psi^{i-1}$ est nul car $\operatorname{Hom}(\Phi^i,\Psi^{i-1})=0$, donc l'hom \deleted{$\unsure{\Psi^i\longrightarrow\Psi|\Psi^i\longrightarrow\Psi^i}$} \inserted{$\Theta=u\alpha{D}:$} $\Psi^{i-1}\mathrel{-}\Psi^{i-1}$ composé avec $n\,\mathrm{id}$ est nul, i.e. $n\Theta=0$, or $\operatorname{Hom}(\Psi^{i-1},\Psi^{i-1})\simeq k$, donc $\Theta=0$, donc $\alpha$ respecte la structure d'extension, et on termine comme plus haut pour conclure $\alpha=0$. Donc en résumé

\[
\operatorname{Hom}(\Phi(i,n),\Phi(j,m))\simeq
\begin{cases}
0\ \text{si } j\unsure{=|\ne}i,\,i+1\\
k\ \text{si }j=i+1\\
\quad(\text{base le composé d'hom can.}\\
\quad\Phi(i,n)\mathrel{-}\Psi^i\longrightarrow\Phi(i+1,m)\unsure{)}\\
k\ \text{si }j=i\ (\text{un hom }\deleted{\text{\unsure{respecte|respectant}}}\ \text{la struct. d'ext.}
\end{cases}
\]

et étant déterminé par ce qu'il induit sur $\Psi^i$ et $\Psi^{i-1}$, qui doit être de la forme $t\,m'\cdot\mathrm{id}_{\Psi^i}$ et $t\,n'\cdot\mathrm{id}_{\Psi^{i}}$, où $n'=\frac n{(m,n)}$, $m'=\frac m{(m,n)}$ (il y a une base can. une fois qu'on a choisi le pgcd $(m,n)$ de $m$ et $n$).

% archive sheet number: 24
